\documentclass[12pt]{article}
\usepackage[utf8]{inputenc}
\usepackage[T1]{fontenc}
\usepackage{hyperref}
\title{Phuie Language Resources for K-12 Curriculum Development}
\author{Compiled by Claude Sonnet 5.5 (claude-sonnet-5-5) for Andr\'as Kornai}
\date{}
\begin{document}
\maketitle

\section{Introduction}
Phuie (ISO 639-3 \texttt{pug}, Glottocode phui1238), also known as Puguli, Pwo, Pwie and Pwien, is a Gur (Niger-Congo, Atlantic-Congo) language of southwestern Burkina Faso, with about 14,000 speakers by one estimate (Joshua Project gives 29,000 for the people group). Speakers call it Phuien and themselves Phuo. It is written in a Latin-based alphabet and is reported as vigorous. The results show a Scripture Earth page and a Webonary site, which suggest some local documentation, but I only saw their landing pages and cannot confirm their contents.

\section{Primary Educational Resources}
No primers or school textbooks were found in searches.

\section{Linguistic and Reference Materials}
\begin{itemize}
\item Puguli language, Wikipedia (English): overview, classification, features such as a base-20 numeral system and two tones. \url{https://en.wikipedia.org/wiki/Puguli_language}
\item Phuie, Hauts-Bassins, Burkina Faso: numerals project data page (MPI-EVA). \url{https://lingweb.eva.mpg.de/channumerals/Phuie.htm}
\item Phuien Webonary site, introduction page (a lexical or documentation site; contents beyond this page not examined). \url{https://www.webonary.work/phuien/overview/introduction/?lang=en}
\end{itemize}

\section{Religious and Community Resources}
\begin{itemize}
\item Scripture Earth page for Phuie (listing of scripture resources; contents not examined). \url{https://www.scriptureearth.org/00eng.php?iso=pug}
\item Pwo in Burkina Faso, Joshua Project: people-group profile. \url{https://joshuaproject.net/people_groups/14489}
\item Phuo Worldview, Burkina Faso (2019 June 3), PDF on the Phuien Webonary site. \url{https://www.webonary.work/phuien/files/2019-June-3-Phuo-Worldview-Burkina-Faso.pdf?lang=en}
\end{itemize}

\section{Digital and Audio Resources}
\begin{itemize}
\item Wikitongues page for Phuie (language profile; EGIDS 6a, vigorous). \url{https://wikitongues.org/languages/pug/}
\end{itemize}

\section{Curriculum Development Context}
The language is vigorous and not known to be taught in schools. The Webonary documentation and the Scripture Earth page are the leads most worth following up with the people behind them, along with an existing Latin-script orthography, to check what dictionary, text and audio material already exists. Note that the wider Puguli/Pwo naming can cause confusion with other groups, so verify identity before reusing any material.

\end{document}